% Autorska priprava k Realisticky.cz 020310–020311; ulohy jsou autorske a obmenene.
% Zamerne nenahrazuje vyklad presnymi kopiemi cizich textu/prikladu.
\documentclass[11pt,a4paper]{article}
\input{../../sablony/pripravy/priprava.sty}
\usepackage{tikz}
\nastavPripravu{1. ročník}{Rovnice a nerovnice}{80}{Soustavy: sčítací metoda}
\begin{document}
\blok{Návaznost a motivace}
Dosazování umíme. Nyní získáme rovnici o jedné neznámé jinak: sečteme nebo odečteme rovnice tak, aby jeden člen vypadl.
\blok{Princip a odvození}
\textbf{Sčítací metoda:} Rovnice můžeme násobit nenulovými čísly, potom je sečíst nebo odečíst. Cílem je odstranit jednu neznámou. Po určení jedné hodnoty dopočítáme druhou a provedeme zkoušku.
\blok{Řešené příklady: od jednoduchého ke složitějšímu}
\begin{prikladReseniBox}[title={Příklad 1}]
Vyřešte $\begin{cases}x+y=7,\\x-y=1.\end{cases}$
\tcblower
\begin{align*}
2x&=8 &&\text{(součet rovnic)} \\
 x&=4 \\
 4+y&=7 \\
 y&=3
\end{align*}

Řešení $[4;3]$ vyhoví oběma rovnicím.
\end{prikladReseniBox}
\begin{prikladReseniBox}[title={Příklad 2}]
Vyřešte $\begin{cases}2x+3y=12,\\3x-2y=5.\end{cases}$
\tcblower
\begin{align*}
2(2x+3y)&=24 \\
 3(3x-2y)&=15 \\
 13x&=39 \\
 x&=3,\quad y=2
\end{align*}

Druhá upravená rovnice je $9x-6y=15$; po součtu s $4x+6y=24$ vypadne $y$.
\end{prikladReseniBox}
\begin{prikladReseniBox}[title={Příklad 3 – sbírka č. 30b}]
Vyřešte sčítací metodou $\begin{cases}3x+2y=20,\\2x+3y=20.\end{cases}$
\tcblower
\begin{align*}
3(3x+2y)&=60,\\
-2(2x+3y)&=-40,\\
5x&=20,\\
x&=4,\\
3\cdot4+2y&=20,\\
y&=4.
\end{align*}
Výsledná dvojice $[4;4]$ splňuje obě původní rovnice.
\end{prikladReseniBox}

\poznamka{Při násobení rovnice smí násobek změnit všechna tři čísla v řádku; typická chyba je násobit pouze člen s neznámou.}
\blok{Samostatné procvičování}
\begin{zadaniReseniBox}[title={Úloha 1}]
$\begin{cases}x+y=9,\\x-y=3.\end{cases}$
\tcblower
\begin{align*}
2x&=12 \\
 x&=6 \\
 y&=3
\end{align*}
\end{zadaniReseniBox}
\begin{zadaniReseniBox}[title={Úloha 2}]
$\begin{cases}2x+y=8,\\2x-y=0.\end{cases}$
\tcblower
\begin{align*}
4x&=8 \\
 x&=2 \\
 y&=4
\end{align*}
\end{zadaniReseniBox}
\begin{zadaniReseniBox}[title={Úloha 3}]
$\begin{cases}3x+2y=16,\\x-y=2.\end{cases}$
\tcblower
\begin{align*}
3x+2y&=16 \\
 2x-2y&=4 \\
 5x&=20\Rightarrow x=4,\ y=2
\end{align*}
\end{zadaniReseniBox}
\blok{Volitelná rezerva}
\begin{zadaniReseniBox}[title={Úloha 1}]
Vyřešte sčítáním $\begin{cases}2x+5y=19,\\3x-5y=6.\end{cases}$
\tcblower
\begin{align*}
5x&=25 \\
 x&=5 \\
 10+5y&=19\Rightarrow y=\frac95
\end{align*}

Zkouška v druhé rovnici: $15-9=6$.
\end{zadaniReseniBox}
\begin{zadaniReseniBox}[title={Úloha ze sbírky č. 30d – pozor, součiny se vyruší}]
Řešte $\begin{cases}(x+4)(y-2)=(x-2)(y+13),\\(x-1)(y-3)=(x+2)(y-5).\end{cases}$
\tcblower
\begin{align*}
xy-2x+4y-8&=xy+13x-2y-26,\\
-15x+6y&=-18,\\
5x-2y&=6,\\[2pt]
xy-3x-y+3&=xy-5x+2y-10,\\
2x-3y&=-13.
\end{align*}
\begin{align*}
3(5x-2y)&=18,\\
-2(2x-3y)&=26,\\
11x&=44,\\
x&=4,\qquad y=7.
\end{align*}
Součiny $xy$ se po úpravě vyruší; výsledkem je lineární soustava. $K=\{[4;7]\}$.
\end{zadaniReseniBox}

\blok{Domácí úkol}
\begin{zadaniReseniBox}[title={Úloha 1}]
$\begin{cases}2x+y=10,\\x-y=2.\end{cases}$
\tcblower
\begin{align*}
3x&=12 \\
 x&=4 \\
 y&=2
\end{align*}
\end{zadaniReseniBox}
\begin{zadaniReseniBox}[title={Úloha 2}]
$\begin{cases}3x+2y=11,\\3x-2y=7.\end{cases}$
\tcblower
\begin{align*}
6x&=18 \\
 x&=3 \\
 y&=1
\end{align*}
\end{zadaniReseniBox}

\poznamka{Orientační časové rozvržení: 5 + 12 + 12 + 10 + 6 = 45 min. Přednost má pochopení jádra; rezerva není součástí časového součtu. Domácí úlohy jsou určeny na samostatnou práci, řešení zůstávají učiteli.}
\end{document}
